% part: intuitionistic-logic
% chapter: soundness-completeness
% section: canonical-model

\documentclass[../../../include/open-logic-section]{subfiles}

\begin{document}

\olfileid{int}{sc}{mod}

\olsection{کانوني مدل}

زموږ په مدل کښې نړۍ به د طبيعي عددونو متناهي لړۍ~$\sigma$
وي، يعنې $\sigma \in \Nat^*$۔ پام وکړئ چې $\Nat^*$ په استقرايي
ډول داسې تعريفېږي:
\begin{enumerate}
\item $\emptyseq \in \Nat^*$۔
\item که $\sigma \in \Nat^*$ او $n \in \Nat$ وي، نو $\sigma.n \in
  \Nat^*$ وي (دلته $\sigma.n$ د $\sigma \concat \tuple{n}$ معنا
  لري او $\sigma \concat \sigma'$ د $\sigma$ او $\sigma'$ يوځاے
  کول دي)۔
\item له دې پرته بل څۀ په $\Nat^*$ کښې نشته۔
\end{enumerate}
نو د استقرايي تعريفونو دپاره $\Nat^*$ کارولے شو۔

$\tuple{!B_0, !C_0}$، $\tuple{!B_1, !C_1}$، \dots، دې د
!!{formula}s د ټولو جوړو يوه شمېرنه وي۔ د !!{formula}s د يوه
سټ~$\Delta$ په لرلو سره، $\Delta(\sigma)$ په استقرايي ډول
داسې تعريفوو:
\begin{enumerate}
\item $\Delta(\emptyseq) = \Delta$
\item $\Delta(\sigma.n) = {}$
  \[
  \begin{cases}
    (\Delta(\sigma) \cup \{!B_n\})^* &
    \text{که $\Delta(\sigma) \cup \{!B_n\} \Proves/ !C_n$ وي} \\
    \Delta(\sigma) & \text{که داسې نۀ وي}
  \end{cases}
  \]
\end{enumerate}
دلته له $(\Delta(\sigma) \cup \{!B_n\})^*$ څخه د
!!{formula}s هغه اولنی سټ مراد دے چې د
\olref[lin]{lem:lindenbaum} له مخې، د $\Delta(\sigma) \cup
\{!B_n\}$ پر سټ او !!{formula}~$!C_n$ د لمې په تطبيقولو شته۔
پام وکړئ چې د دې تعريف له مخې، که $\Delta(\sigma) \cup
\{!B_n\} \Proves/ !C_n$ وي، نو $\Delta(\sigma.n)
\Proves !B_n$ او $\Delta(\sigma.n) \Proves/ !C_n$ وي۔ دا هم
پام وکړئ چې د هر~$n$ دپاره $\Delta(\sigma) \subseteq
\Delta(\sigma.n)$ وي۔ که $\Delta$ اولنی وي، نو د هر~$\sigma$
دپاره $\Delta(\sigma)$ هم اولنی وي۔

\begin{defn}\ollabel{defn:canonical-model}
  فرض کړئ $\Delta$ اولنی دے۔ نو د~$\Delta$ دپاره
  \emph{کانوني مدل} $\mModel{M(\Delta)}$ داسې تعريفېږي:
  \begin{enumerate}
  \item $W = \Nat^*$، د طبيعي عددونو د متناهي لړيو سټ۔
  \item $R$ داسې جزئي ترتيب دے چې $R\sigma\sigma'$ هله او
    يوازې هله وي چې $\sigma$ د~$\sigma'$ لومړنۍ برخه وي
    (يعنې د کومې لړۍ~$\sigma''$ دپاره $\sigma' =
    \sigma \concat \sigma''$ وي)۔
  \item $V(p) = \Setabs{\sigma}{p \in \Delta(\sigma)}$۔
  \end{enumerate}
\end{defn}

د پيلنۍ برخې اړيکه~$R$ جزئي ترتيب دے۔ هره لړۍ خپله پيلنۍ
برخه ده؛ د پيلنۍ برخې پيلنۍ برخه هم پيلنۍ برخه ده؛ او که
دوې لړۍ د يو بل پيلنۍ برخې وي، يو شان اوږدوالے او يو شان غړي
لري۔ د~$V$ د يورخيزې زياتېدنې شرط هم پوره دے۔ ځکه چې $\Delta(\sigma)
\subseteq \Delta(\sigma.n)$ وي، نو له~$\sigma$ وروسته د ورزياتې
شوې برخې پر اوږدوالي په استقرا ښيو چې هر کله $R\sigma\sigma'$ وي،
$\Delta(\sigma)
\subseteq \Delta(\sigma')$ راځي۔ که ورزياته
شوې برخه تشه وي، سټونه برابر دي۔ هر بل زياتېدلے غړے د سټ
شمول ساتي۔ نو هر اټومي متغير چې په لومړي سټ کښې وي، په
دويم کښې هم وي۔

\end{document}
